\documentclass[11pt,a4paper]{article}
\usepackage[utf8]{inputenc}
\usepackage[T1]{fontenc}
\usepackage[french]{babel}
\usepackage{lmodern}
\usepackage{geometry}
\geometry{hmargin=2.5cm,vmargin=2.5cm}
\usepackage{xcolor}
\usepackage{amsmath,amssymb}
\usepackage{graphicx}
\usepackage{booktabs}
\usepackage{tcolorbox}
\tcbuselibrary{skins}
\usepackage{titlesec}
\usepackage{fancyhdr}
\usepackage{hyperref}

% Définition des couleurs du thème Lacuna
\definecolor{deepblue}{HTML}{1E88E5}
\definecolor{deepred}{HTML}{E53935}
\definecolor{gamebg}{HTML}{14161C}
\definecolor{cardbg}{HTML}{19171D}
\definecolor{cyanflower}{HTML}{00D2D2}
\definecolor{pinkflower}{HTML}{DC64B4}
\definecolor{purpleflower}{HTML}{8C3CC8}
\definecolor{greenflower}{HTML}{32B432}
\definecolor{orangeflower}{HTML}{FFA500}
\definecolor{redflower}{HTML}{DC3232}
\definecolor{blueflower}{HTML}{3282DC}

\hypersetup{
    colorlinks=true,
    linkcolor=deepblue,
    filecolor=magenta,      
    urlcolor=deepblue,
    pdftitle={Dossier de Validation IHM - Jeu Lacuna},
    pdfpagemode=FullScreen,
}

% Configuration de tcolorbox pour les dialogues de tests
\tcbset{
    dialoguebox/.style={
        colback=white,
        colframe=gray!30,
        arc=3mm,
        left=4mm,
        right=4mm,
        top=3mm,
        bottom=3mm,
        boxrule=0.5pt,
        fonttitle=\bfseries\color{black},
        colbacktitle=gray!10,
        attach boxed title to top left={yshift=-2mm, xshift=4mm},
        boxed title style={boxrule=0.5pt, colframe=gray!30, arc=1mm}
    },
    gamerbox/.style={
        dialoguebox,
        colframe=deepred!40,
        colbacktitle=deepred!10,
        fonttitle=\bfseries\color{deepred}
    },
    nongamerbox/.style={
        dialoguebox,
        colframe=deepblue!40,
        colbacktitle=deepblue!10,
        fonttitle=\bfseries\color{deepblue}
    },
    seniorbox/.style={
        dialoguebox,
        colframe=orangeflower!40,
        colbacktitle=orangeflower!10,
        fonttitle=\bfseries\color{orangeflower!80!black}
    }
}

% Styles de sections
\titleformat{\section}{\large\bfseries\color{gamebg!80!blue}}{\thesection}{1em}{}[{\titlerule[0.5pt]}]
\titleformat{\subsection}{\normalfont\bfseries\color{deepblue!80!black}}{\thesubsection}{1em}{}

% En-têtes et pieds de page
\pagestyle{fancy}
\fancyhf{}
\lhead{\color{gray}Dossier de Validation IHM -- Jeu Lacuna}
\rhead{\color{gray}PROG6 -- S6}
\cfoot{\thepage}
\renewcommand{\headrulewidth}{0.4pt}
\renewcommand{\footrulewidth}{0.4pt}

\begin{document}

% --- Page de Garde ---
\begin{titlepage}
    \centering
    \vspace*{1.5cm}
    {\large\bfseries Université Grenoble Alpes (UGA)}\\[0.2cm]
    {\large UFR IM2AG -- Licence 3 Informatique}\\[0.5cm]
    {\large Module : PROG6 -- Semestre 6}\\[2.5cm]
    
    \rule{\linewidth}{1.5pt}\\[0.5cm]
    {\huge\bfseries Dossier de Validation IHM}\\[0.3cm]
    {\Large\itshape Évaluation Ergonomique et Tests Utilisateurs du Jeu Lacuna}\\[0.5cm]
    \rule{\linewidth}{1.5pt}\\[2.5cm]
    
    \begin{minipage}[t]{0.48\textwidth}
        \begin{flushleft} \large
            \emph{Auteurs :}\\
            MADJOUR Amir\\
            FRITAH Wassim\\
            ZOUITEN Ouassim\\
            BOUGUERRA Lakhdar\\
            DJERBOA Mohamed\\
            SAHRAOUI Mohammed
        \end{flushleft}
    \end{minipage}
    \hfill
    \begin{minipage}[t]{0.45\textwidth}
        \begin{flushright} \large
            \emph{Enseignant tuteur :}\\
            Benjamin Wack
        \end{flushright}
    \end{minipage}\\[1.5cm]
    
    \vfill
    {\large \today}
\end{titlepage}

\tableofcontents
\newpage

% --- Section 1 : Introduction et Fonctionnement du Jeu ---
\section{Introduction et Fonctionnement du Jeu}

Le projet \textbf{Lacuna} est l'adaptation logicielle en Java Swing d'un jeu de société abstrait et géométrique. Avant d'aborder les aspects ergonomiques de notre interface, il convient de comprendre les règles du jeu physique.

\subsection{But du Jeu et Matériel}
Le jeu se déroule sur un tapis de jeu circulaire sur lequel sont réparties de manière pseudo-aléatoire \textbf{49 fleurs} appartenant à \textbf{7 couleurs (espèces) différentes} (soit 7 fleurs par couleur). Deux joueurs s'affrontent avec pour objectif de remporter la majorité des fleurs dans au moins 4 des 7 couleurs disponibles.

\subsection{Déroulement d'une Partie}
Une partie se structure en deux phases successives :
\begin{enumerate}
    \item \textbf{La phase de placement et de capture directe :}
    Chaque joueur dispose de \textbf{6 pions}. À tour de rôle, les joueurs placent un de leurs pions sur le plateau. Pour ce faire, le pion doit être posé sur un segment de droite imaginaire reliant deux fleurs libres de la même couleur. Le joueur capture alors instantanément ces deux fleurs, qui sont retirées du plateau.
    Deux contraintes fondamentales s'appliquent :
    \begin{itemize}
        \item Le segment de droite reliant les deux fleurs de couleur identique ne doit être obstrué par aucun autre élément (ni pion déjà posé, ni autre fleur présente sur le plateau).
        \item Le pion ne peut pas être placé trop près d'une fleur (respect d'une distance physique minimale).
    \end{itemize}
    De plus, lors du démarrage de la partie, un tirage au sort désigne une fleur de départ qui est attribuée gratuitement au premier joueur.
    
    \item \textbf{La phase de résolution par proximité (Le "Hold-up") :}
    Une fois que les 12 pions (6 par joueur) ont été posés sur le plateau, la phase de placement prend fin. Les fleurs restantes sur le tapis (qui n'ont pas pu être capturées directement) sont alors attribuées automatiquement au joueur dont l'un des pions est géométriquement le plus proche de chacune d'elles (résolution par diagramme de Voronoï).
\end{enumerate}

\subsection{Détermination du Vainqueur}
Pour chacune des 7 couleurs, on compare le nombre de fleurs collectées par chaque joueur. Celui qui possède le plus grand nombre de fleurs d'une couleur (soit au moins 4 fleurs sur les 7) remporte la majorité pour cette couleur. Le joueur qui obtient au moins \textbf{4 majorités} gagne la partie.

\subsection{Problématiques et Défis de Conception Rencontrés}
L'adaptation de ce jeu de société physique en version numérique soulève plusieurs questions et problèmes ergonomiques majeurs que nous avons dû résoudre :
\begin{itemize}
    \item \textbf{Comment représenter les fleurs ?} Les 49 fleurs de 7 espèces (couleurs) différentes doivent être immédiatement identifiables sur le plateau pour que les joueurs puissent repérer les alignements stratégiques.
    \item \textbf{Comment représenter les pions ?} Les pions placés par les joueurs doivent se distinguer clairement des pions-fleurs environnants pour marquer le contrôle de l'espace sur le plateau.
\item \textbf{En cas d'intersection, comment le joueur peut-il choisir la ligne qu'il veut ?} Si plusieurs alignements valides se croisent sous le pointeur, le système doit permettre de sélectionner précisément celle qu'il souhaite capturer.
    \item \textbf{Comment représenter le tour du joueur ?} Informer clairement et instantanément le joueur actif que c'est à lui de jouer, tout en maintenant sa concentration sur la zone de jeu.
    \item \textbf{Comment représenter les compteurs de fleurs ?} Afficher les scores en temps réel de manière à ce que la course aux majorités (l'objectif de collecter au moins 4 fleurs sur 7 par couleur) soit immédiatement décryptable.
\end{itemize}

\newpage

% --- Section 2 : Conception Itérative de l'Interface ---
\section{Conception Itérative de l'Interface}

\subsection{Comment représenter les fleurs ?}

La toute première question que nous avons posée était visuelle : le plateau de jeu est peuplé de 49 fleurs réparties en 7 espèces (couleurs). Comment les rendre immédiatement distinguishables les unes des autres, sans grille ni légende, sur un écran ?

Nous avons opté pour une palette de 7 teintes vives (Cyan, Rose, Violet, Vert, Orange, Rouge, Bleu) sur un fond très sombre (\texttt{\#14161C}). Le contraste maximal devait permettre au regard d'identifier instantanément les regroupements de couleurs. Nous avons implémenté cette charte et joué plusieurs parties en interne.

Les tests internes ont confirmé la lisibilité générale. Nous avons ensuite soumis l'interface à nos testeurs. Voici ce qu'ils ont observé :

\begin{tcolorbox}[nongamerbox, title={Djaber -- 23 ans, étudiant en design graphique}]
    \textbf{Q. Comment juges-tu la lisibilité des 7 couleurs de fleurs sur le plateau ?}\\
    \textbf{Djaber :} L'harmonie des couleurs est bien choisie. Sur le fond sombre, le cyan, le rose, le vert ou le jaune ressortent parfaitement. Cependant, j'ai eu un léger doute au début entre les fleurs rouges et les fleurs roses lorsque le plateau était très chargé en pions. Peut-être qu'il faudrait donner des formes légèrement différentes aux pétales selon les couleurs pour aider les personnes qui ont du mal à distinguer les teintes proches.
\end{tcolorbox}

Ce retour a été confirmé par le panel senior lors de parties en fin de plateau dense. Le risque de confusion rouge/rose a été identifié comme une limite réelle. En réponse, nous avons acté l'ajout d'une option d'accessibilité future permettant de varier la forme géométrique des pétales selon l'espèce de la fleur, afin d'offrir un second canal d'identification au-delà de la couleur.

\subsection{Comment représenter les pions des joueurs ?}

Une fois les fleurs visuellement définies, la question suivante s'est posée d'elle-même : comment représenter les pions que chaque joueur place sur le plateau, de sorte qu'ils se distinguent clairement des fleurs tout en identifiant immédiatement quel joueur les a posés ?

Nous avons choisi des pions circulaires de grande taille avec un design métallique, colorés aux couleurs des joueurs : Rouge Corail intense (\texttt{\#E53935}) pour le Joueur 1, Bleu Électrique (\texttt{\#1E88E5}) pour le Joueur 2. Les tests internes ont validé immédiatement la différenciation. Les testeurs ont également confirmé cette clarté :

\begin{tcolorbox}[gamerbox, title={Nadir -- 22 ans, joueur de jeux de cartes en ligne}]
    \textbf{Q. Comment qualifies-tu l'esthétique générale de l'interface ?}\\
    \textbf{Nadir :} C'est très moderne. Le thème sombre avec les couleurs de néon pour les fleurs donne un côté haut de gamme qui change des IHM Java Swing habituelles qui font très ``années 90''. On sent qu'il y a eu un vrai travail sur les ombres portées et les halos.
\end{tcolorbox}

La représentation des pions a été validée sans nécessiter de réitération. Les pions restent très lisibles à toutes les phases de la partie.

\subsection{Comment le joueur peut-il placer un pion et capturer une paire de fleurs ?}

C'est la question centrale d'interaction du jeu : pour poser un pion, le joueur doit viser un segment imaginaire reliant deux fleurs libres de même couleur, sans aucune grille pour l'y aider. Comment matérialiser ce geste géométrique de façon naturelle ?

Notre solution a été de masquer le curseur système et de le remplacer par un \textbf{pion fantôme} translucide de la couleur du joueur actif. Dès que ce pion fantôme survole une zone alignée avec deux fleurs compatibles, une ligne de visée colorée apparaît dynamiquement. En cliquant, les deux fleurs s'envolent vers la grille de score en une animation courbe.

Nous avons testé cette solution en interne et trouvé l'effet immersif et intuitif. Nous avons ensuite demandé à nos testeurs de jouer :

\begin{tcolorbox}[nongamerbox, title={Marouan -- 25 ans, étudiant en lettres}]
    \textbf{Q. Comment as-tu trouvé le système de placement des pions ?}\\
    \textbf{Marouan :} Très naturel. On déplace la souris, et dès qu'on passe entre deux fleurs de la même couleur, le jeu nous propose un trait coloré. On clique et le pion se pose. C'est magique et très simple.
\end{tcolorbox}

\begin{tcolorbox}[seniorbox, title={Le père d'Amir -- 68 ans, ancien cadre administratif}]
    \textbf{Q. Le déplacement du pion fantôme à l'écran vous a-t-il paru stable ?}\\
    \textbf{Le père d'Amir :} Au début, j'avais un peu de mal à viser précisément le milieu de la ligne entre deux fleurs, car ma main tremble un tout petit peu sur la souris. Mais j'ai remarqué que le jeu ``aimantait'' le pion dès que je m'approchais de la ligne. C'est une excellente idée ! Ça compense mon manque de précision et évite de rater son coup par accident.
\end{tcolorbox}

Ce dernier retour a mis en lumière un problème de précision motrice pour le panel senior. Le pion fantôme seul ne suffisait pas --- l'alignement géométrique est trop fin pour être ciblé avec précision à la souris. Nous avons alors intégré une \textbf{aimantation adaptative} : dès que le curseur entre dans un rayon de proximité d'un alignement valide, le pion se cale automatiquement dessus. Cette itération a été re-testée et validée par l'ensemble des testeurs.

\begin{tcolorbox}[nongamerbox, title={Khaled -- 24 ans, alternant en commerce}]
    \textbf{Q. Le système de retour en arrière (Undo/Redo) t'a-t-il été utile ?}\\
    \textbf{Khaled :} Oui, pendant ma première partie contre l'IA, j'ai glissé avec la souris et j'ai cliqué à côté de la ligne que je voulais. J'ai vu le bouton ``Annuler'' en bas, j'ai cliqué dessus et mon pion est revenu dans ma main. C'est super rassurant de savoir qu'on n'est pas puni pour un faux mouvement.
\end{tcolorbox}

Ce retour a conduit à l'ajout des raccourcis clavier standards \texttt{Ctrl+Z} / \texttt{Ctrl+Y} pour l'annulation et le rétablissement, rendant cette fonctionnalité accessible aux profils plus technophiles sans modifier l'interface principale.

\subsection{En cas d'intersection, comment le joueur peut-il choisir la ligne qu'il veut ?}

En introduisant le pion fantôme et les lignes de visée, nous avons immédiatement rencontré un nouveau problème pendant nos propres tests de jeu : lorsque plusieurs alignements valides se croisent sous le curseur, le système doit en choisir un. Lequel afficher ? Lequel permettre au joueur de sélectionner ?

Notre première implémentation laissait le moteur sélectionner arbitrairement le premier alignement détecté. Lors de nos tests internes, cela s'est avéré injouable dans les zones denses du plateau : la ligne affichée n'était pas celle souhaitée, sans moyen de la changer. Nous avons retravaillé la solution en permettant à l'utilisateur de faire défiler la \textbf{molette de la souris} pour cycler entre les lignes candidates. La ligne active s'affiche en trait plein, les autres en pointillés légers.

\begin{tcolorbox}[gamerbox, title={Radhi -- 23 ans, développeur amateur et joueur de jeux de plateau}]
    \textbf{Q. L'ergonomie de la molette pour choisir entre plusieurs lignes de fleurs t'a-t-elle convaincu ?}\\
    \textbf{Radhi :} C'est une excellente idée d'interaction. J'avais peur que ce soit brouillon quand trois lignes de fleurs roses se croisent au même endroit, mais le petit popup ``Molette : Choisir la ligne'' qui apparaît à côté du curseur est parfait. Tu scrolles et la ligne active s'épaissit tandis que les autres passent en pointillés légers. C'est une très bonne solution au problème des chevauchements.
\end{tcolorbox}

Cependant, un autre retour est venu nuancer cette validation :

\begin{tcolorbox}[seniorbox, title={La mère de Wassim -- 70 ans, retraitée de la fonction publique}]
    \textbf{Q. Est-ce que l'utilisation de la molette pour choisir la ligne de fleurs a été facile pour vous ?}\\
    \textbf{La mère de Wassim :} J'ai eu un peu de mal à comprendre ce geste au début, car je n'utilise jamais la molette de ma souris d'habitude. Je cherchais à cliquer sur les fleurs pour choisir. Après une explication, j'ai compris le coup de main, mais pour les personnes de mon âge, un simple bouton ``Suivant'' à l'écran pour changer de ligne aurait peut-être été plus rassurant qu'un mouvement de molette.
\end{tcolorbox}

La molette est validée pour les profils habitués, mais nous avons pris en compte ce retour pour le définir comme piste d'évolution future : ajouter une alternative visuelle cliquable pour les utilisateurs peu familiers avec le défilement matériel.

\subsection{Comment représenter le tour du joueur ?}

Avec l'interaction de base fonctionnelle, la question du déroulement de la partie s'est posée : comment le joueur sait-il instantanément que c'est à son tour de jouer, sans avoir à chercher l'information à l'écran ?

Notre première approche était une simple indication textuelle dans le panneau du joueur actif. Lors des sessions de jeu en interne, nous avons constaté que des joueurs rataient leur tour ou hésitaient, l'indication étant trop discrète. Nous avons alors remplacé cela par un \textbf{halo lumineux pulsé} autour du panneau du joueur actif, dans sa couleur (rouge ou bleu).

\begin{tcolorbox}[gamerbox, title={Radhi -- 23 ans, développeur amateur et joueur de jeux de plateau}]
    \textbf{Q. As-tu testé le jeu en réseau ? L'IHM gère-t-elle bien l'attente ?}\\
    \textbf{Radhi :} Oui, j'ai fait une partie en ligne avec Nadir. La boîte de dialogue de déconnexion et l'écran d'attente d'action de l'autre joueur sont clairs. On sait toujours si c'est notre tour ou si l'adversaire réfléchit, notamment grâce au halo lumineux pulsé autour du nom du joueur actif.
\end{tcolorbox}

Le halo pulsé a été validé sans nécessiter de nouvelle itération, par tous les profils testeurs.

\subsection{Comment représenter les compteurs de fleurs capturées ?}

La dernière grande question de conception concernait l'affichage du score : comment permettre à chaque joueur de savoir, en temps réel et d'un seul coup d'œil, combien de fleurs il a capturées par couleur et s'il est en train de remporter une majorité ?

Nous avons opté pour des panneaux latéraux avec une table verticale associant le sprite coloré de chaque fleur à son compteur. Après nos premières sessions internes, la lisibilité nous a semblé bonne. Mais les testeurs du panel senior ont mis en lumière une limite :

\begin{tcolorbox}[seniorbox, title={La mère d'Amir -- 66 ans, retraitée de l'enseignement}]
    \textbf{Q. Avez-vous réussi à lire facilement les textes et les scores à l'écran ?}\\
    \textbf{La mère d'Amir :} Le contraste est très bon pour les fleurs et le plateau, car l'écran est bien sombre et les éléments sont lumineux. En revanche, j'ai eu quelques difficultés à lire les petits chiffres des scores dans les panneaux latéraux. Mes lunettes ne suffisaient pas toujours, j'aurais aimé pouvoir agrandir la taille de ces textes ou avoir une option pour basculer en haute visibilité.
\end{tcolorbox}

En revanche, pour les autres profils, la lisibilité a été validée :

\begin{tcolorbox}[gamerbox, title={Nadir -- 22 ans, joueur de jeux de cartes en ligne}]
    \textbf{Q. Le tableau de scores latéral est-il facile à lire en pleine partie ?}\\
    \textbf{Nadir :} Oui, la disposition verticale avec les icônes de fleurs et les chiffres est très lisible. On voit tout de suite qui possède la majorité sur quelle couleur, ce qui est crucial pour la stratégie de fin de partie.
\end{tcolorbox}

Suite à ces retours, la structure des compteurs a été conservée, et nous avons acté l'ajout d'un mode accessibilité senior comme évolution future, permettant de doubler la taille des caractères et d'activer un affichage à fort contraste.

\newpage

% --- Section 3 : Validation par les Testeurs ---
\section{Validation par les Testeurs}

Pour valider l'ensemble de nos choix d'interface, nous avons organisé une campagne de tests auprès de 12 personnes réparties en 3 profils distincts.

\begin{table}[ht]
    \centering
    \caption{Groupes de testeurs}
    \vspace{0.2cm}
    \begin{tabular}{llp{7.5cm}l}
        \toprule
        \textbf{Groupe} & \textbf{Public ciblé} & \textbf{Objectif principal} & \textbf{Âge moyen} \\
        \midrule
        \textbf{Groupe 1} & Amis Gamers & Efficacité, rapidité, retour d'information poussé & 22 ans \\
        \textbf{Groupe 2} & Amis Non-Gamers & Facilité d'apprentissage, clarté, intuitivité & 24 ans \\
        \textbf{Groupe 3} & Personnes Âgées & Accessibilité, taille de police, tolérance aux erreurs & 71 ans \\
        \bottomrule
    \end{tabular}
    \label{tab:groupes}
\end{table}

\subsection{Groupe 1 : Nos amis Gamers}

\begin{tcolorbox}[gamerbox, title={Mohammed -- 21 ans, joueur d'esport compétitif}]
    \textbf{Q. Comment évalues-tu l'intuitivité du système de visée ?}\\
    \textbf{Mohammed :} Franchement, c'est super propre. Le fait que le pointeur de la souris disparaisse pour laisser place au pion fantôme avec les lignes d'alignement qui apparaissent dès qu'on passe sur le plateau, ça permet d'enchaîner super vite. Le feedforward est impeccable, pas besoin de réfléchir pour comprendre si le coup est valide ou pas.

    \textbf{Q. Qu'as-tu pensé de l'écran de Voronoï lors du Hold-up final ?}\\
    \textbf{Mohammed :} C'est le gros point fort technique de votre interface. En général, dans ce genre de jeu, tu dois calculer les distances à l'œil nu et c'est chiant. Là, les polygones de Voronoï s'affichent avec les lignes élastiques colorées. C'est instantané, je sais exactement pourquoi j'ai gagné ou perdu telle fleur.

    \textbf{Q. As-tu des suggestions d'amélioration ?}\\
    \textbf{Mohammed :} Pour le rythme, j'aimerais bien pouvoir zapper l'animation d'intro du cylindre au début en cliquant. Et ce serait cool d'ajouter des raccourcis clavier comme \texttt{Ctrl+Z} et \texttt{Ctrl+Y} pour l'annulation/rétablissement.
\end{tcolorbox}

\begin{tcolorbox}[gamerbox, title={Ilyes -- 22 ans, amateur de jeux de stratégie}]
    \textbf{Q. As-tu rencontré des difficultés à comprendre les règles d'obstruction via l'IHM ?}\\
    \textbf{Ilyes :} Non, c'est très clair. La ligne devient rouge et un toast discret s'affiche en haut de l'écran avec un message du type ``La ligne traverse une autre pièce'' ou ``Le pion est trop près''. On sait tout de suite ce qui coince.

    \textbf{Q. Que penses-tu du rythme global des animations de capture ?}\\
    \textbf{Ilyes :} C'est visuellement très flatteur, les fleurs tournoient bien vers le tableau de score. Mais quand on commence à bien connaître le jeu et qu'on veut enchaîner, la transition d'envol des fleurs est peut-être un poil lente. Ce serait bien d'avoir une option pour passer en mode ``animations rapides''.
\end{tcolorbox}

\begin{tcolorbox}[gamerbox, title={Radhi -- 23 ans, développeur amateur et joueur de jeux de plateau}]
    \textbf{Q. L'ergonomie de la molette pour choisir entre plusieurs lignes de fleurs t'a-t-elle convaincu ?}\\
    \textbf{Radhi :} C'est une excellente idée d'interaction. J'avais peur que ce soit brouillon quand trois lignes de fleurs roses se croisent au même endroit, mais le petit popup ``Molette : Choisir la ligne'' qui apparaît à côté du curseur est parfait. Tu scrolles et la ligne active s'épaissit tandis que les autres passent en pointillés légers.

    \textbf{Q. As-tu testé le jeu en réseau ? L'IHM gère-t-elle bien l'attente ?}\\
    \textbf{Radhi :} Oui, j'ai fait une partie en ligne avec Nadir. La boîte de dialogue de déconnexion et l'écran d'attente d'action de l'autre joueur sont clairs. On sait toujours si c'est notre tour ou si l'adversaire réfléchit, notamment grâce au halo lumineux pulsé autour du nom du joueur actif.
\end{tcolorbox}

\begin{tcolorbox}[gamerbox, title={Nadir -- 22 ans, joueur de jeux de cartes en ligne}]
    \textbf{Q. Comment qualifies-tu l'esthétique générale de l'interface ?}\\
    \textbf{Nadir :} C'est très moderne. Le thème sombre avec les couleurs de néon pour les fleurs donne un côté haut de gamme qui change des IHM Java Swing habituelles qui font très ``années 90''. On sent qu'il y a eu un vrai travail sur les ombres portées et les halos.

    \textbf{Q. Le tableau de scores latéral est-il facile à lire en pleine partie ?}\\
    \textbf{Nadir :} Oui, la disposition verticale avec les icônes de fleurs et les chiffres est très lisible. On voit tout de suite qui possède la majorité sur quelle couleur, ce qui est crucial pour la stratégie de fin de partie.
\end{tcolorbox}

\subsection{Groupe 2 : Nos amis Non-Gamers}

\begin{tcolorbox}[nongamerbox, title={Marouan -- 25 ans, étudiant en lettres}]
    \textbf{Q. As-tu utilisé le tutoriel interactif ? T'a-t-il aidé à comprendre le jeu ?}\\
    \textbf{Marouan :} Oui, j'ai lancé le tutoriel dès le début car je ne connaissais pas du tout les règles de Lacuna. C'est super bien fait. Le système d'étapes avec l'écran qui s'assombrit pour mettre en valeur les éléments dont on parle m'a permis de comprendre le jeu en deux minutes.

    \textbf{Q. Comment as-tu trouvé le système de placement des pions ?}\\
    \textbf{Marouan :} Très naturel. On déplace la souris, et dès qu'on passe entre deux fleurs de la même couleur, le jeu nous propose un trait coloré. On clique et le pion se pose. C'est magique et très simple.
\end{tcolorbox}

\begin{tcolorbox}[nongamerbox, title={Djaber -- 23 ans, étudiant en design graphique}]
    \textbf{Q. Comment juges-tu la lisibilité des 7 couleurs de fleurs sur le plateau ?}\\
    \textbf{Djaber :} L'harmonie des couleurs est bien choisie. Sur le fond sombre, le cyan, le rose, le vert ou le jaune ressortent parfaitement. Cependant, j'ai eu un léger doute au début entre les fleurs rouges et les fleurs roses lorsque le plateau était très chargé en pions.

    \textbf{Q. As-tu apprécié les animations visuelles ?}\\
    \textbf{Djaber :} Totalement. L'envol des fleurs vers les compteurs est hyper gratifiant. On comprend tout de suite que la fleur est capturée et où elle va. Ça apporte une vraie satisfaction visuelle.
\end{tcolorbox}

\begin{tcolorbox}[nongamerbox, title={Khaled -- 24 ans, alternant en commerce}]
    \textbf{Q. Le système de retour en arrière (Undo/Redo) t'a-t-il été utile ?}\\
    \textbf{Khaled :} Oui, pendant ma première partie contre l'IA, j'ai glissé avec la souris et j'ai cliqué à côté de la ligne que je voulais. J'ai vu le bouton ``Annuler'' en bas, j'ai cliqué dessus et mon pion est revenu dans ma main. C'est super rassurant de savoir qu'on n'est pas puni pour un faux mouvement.

    \textbf{Q. Qu'as-tu pensé de l'IA du jeu ? L'IHM rend-elle ses coups lisibles ?}\\
    \textbf{Khaled :} L'IA joue assez vite, mais l'IHM affiche un message ``L'ordinateur réfléchit...'' puis anime le coup. Le jeu trace la ligne choisie par l'ordinateur pendant une seconde pour qu'on comprenne ce qu'il vient de faire avant de passer à notre tour. C'est très bien dosé.
\end{tcolorbox}

\begin{tcolorbox}[nongamerbox, title={Sami -- 24 ans, infirmier}]
    \textbf{Q. Est-ce que le jeu t'a paru difficile à prendre en main sur le plan informatique ?}\\
    \textbf{Sami :} Pas du tout. Je ne joue jamais sur ordinateur, mais là, tout se fait à la souris. On n'a pas besoin d'utiliser le clavier. Le fait qu'on ne puisse pas poser un pion à un endroit interdit simplifie beaucoup les choses.

    \textbf{Q. Les fenêtres de dialogue pour saisir les noms ou choisir qui commence étaient-elles claires ?}\\
    \textbf{Sami :} Oui, très simples. Le bouton rouge et le bouton bleu pour choisir le joueur qui commence sont gros et explicites. On clique sur le visage ou le nom et le jeu démarre immédiatement.
\end{tcolorbox}

\subsection{Groupe 3 : Les Personnes Âgées}

\begin{tcolorbox}[seniorbox, title={La mère d'Amir -- 66 ans, retraitée de l'enseignement}]
    \textbf{Q. Avez-vous réussi à lire facilement les textes et les scores à l'écran ?}\\
    \textbf{La mère d'Amir :} Le contraste est très bon pour les fleurs et le plateau, car l'écran est bien sombre et les éléments sont lumineux. En revanche, j'ai eu quelques difficultés à lire les petits chiffres des scores dans les panneaux latéraux. Mes lunettes ne suffisaient pas toujours, j'aurais aimé pouvoir agrandir la taille de ces textes.

    \textbf{Q. Les boutons d'interface (Quitter, Rejouer, Annuler) étaient-ils assez grands ?}\\
    \textbf{La mère d'Amir :} Oui, les boutons en bas de l'écran sont bien espacés et larges. On clique dessus sans risquer de se tromper de bouton. C'est bien pensé.
\end{tcolorbox}

\begin{tcolorbox}[seniorbox, title={Le père d'Amir -- 68 ans, ancien cadre administratif}]
    \textbf{Q. Le déplacement du pion fantôme à l'écran vous a-t-il paru stable ?}\\
    \textbf{Le père d'Amir :} Au début, j'avais un peu de mal à viser précisément le milieu de la ligne entre deux fleurs, car ma main tremble un tout petit peu sur la souris. Mais j'ai remarqué que le jeu ``aimantait'' le pion dès que je m'approchais de la ligne. Ça compense mon manque de précision et évite de rater son coup par accident.

    \textbf{Q. Avez-vous compris le principe du calcul des scores à la fin ?}\\
    \textbf{Le père d'Amir :} Oui, grâce aux zones colorées en rouge et bleu qui sont apparues à la fin. On voit tout de suite quelle partie du tapis nous appartient. C'est beaucoup plus clair que si le jeu s'était contenté d'afficher un tableau de chiffres compliqués.
\end{tcolorbox}

\begin{tcolorbox}[seniorbox, title={L'oncle de Wassim -- 74 ans, ancien artisan électricien}]
    \textbf{Q. Comment avez-vous trouvé le rythme visuel du jeu (les animations, le mouvement des pièces) ?}\\
    \textbf{L'oncle de Wassim :} C'est un joli spectacle au début quand toutes les fleurs tombent sur le plateau. Cependant, au milieu de la partie, les mouvements rapides des fleurs qui s'envolent m'ont un peu fatigué les yeux. J'aurais préféré des mouvements plus doux ou plus lents.

    \textbf{Q. Le tutoriel vous a-t-il été utile ?}\\
    \textbf{L'oncle de Wassim :} Oui, le guide pas à pas est indispensable. J'ai pu le lire à mon rythme en cliquant sur ``Continuer'' à chaque étape, ce qui m'a rassuré avant de commencer ma première partie contre la machine.
\end{tcolorbox}

\begin{tcolorbox}[seniorbox, title={La mère de Wassim -- 70 ans, retraitée de la fonction publique}]
    \textbf{Q. Est-ce que l'utilisation de la molette pour choisir la ligne de fleurs a été facile pour vous ?}\\
    \textbf{La mère de Wassim :} J'ai eu un peu de mal à comprendre ce geste au début, car je n'utilise jamais la molette de ma souris d'habitude. Je cherchais à cliquer sur les fleurs pour choisir. Après une explication, j'ai compris le coup de main, mais un simple bouton ``Suivant'' à l'écran aurait peut-être été plus rassurant.

    \textbf{Q. Le fait de jouer sur un fond très sombre vous a-t-il fatigué les yeux ?}\\
    \textbf{La mère de Wassim :} Au contraire, c'est très reposant. Les écrans blancs me font souvent mal aux yeux après quelques minutes. Ici, le fond noir permet de bien se concentrer sur les pions de couleur sans aucune fatigue visuelle.
\end{tcolorbox}

\newpage

% --- Section 4 : Conclusion et Perspectives ---
\section{Conclusion et Perspectives}

La campagne d'évaluation menée auprès des 12 participants a validé la grande majorité de nos choix d'interface. L'assistance visuelle au placement (pion fantôme, aimantation, coloration des lignes) remplit son rôle de feedforward pour tous les publics. Le partitionnement de Voronoï rend la règle géométrique de fin de partie accessible et transparente. La molette de sélection est appréciée des profils technophiles.

Les retours identifient deux évolutions prioritaires à implémenter :
\begin{itemize}
    \item \textbf{Un mode accessibilité senior :} Agrandissement des compteurs de score et différenciation des fleurs par la forme des pétales, pour les personnes malvoyantes ou daltoniennes.
    \item \textbf{Une alternative à la molette :} Un bouton cliquable à l'écran pour changer de ligne candidate, afin d'éliminer la dépendance au défilement matériel pour les utilisateurs peu familiers.
\end{itemize}

\end{document}
